\chapter{引言与环境准备}
\label{chap:intro}

本章用于建立阅读 DBMS\_C 的共同起点。读者在进入具体模块实验之前，需要先知道本项目的教学定位、源码目录如何组织、如何稳定构建和测试，以及本指导书为什么采用“先测试、再写实验内容”的组织方式。

\section{编写目的}
\label{sec:intro-purpose}

本指导书面向计算机专业低年级学生、数据库课程实践和 DBMS\_C 项目阅读者。它的目标不是让读者一次性掌握所有数据库实现细节，而是通过“源码阅读 + 测试运行 + Trace 观察 + 图表理解”的方式，逐步理解一个轻量化 DBMS 原型从 SQL 输入到文件存储访问的大致过程。

各章会围绕真实源码展开：先确认或补充可运行的测试，再把测试目标、运行命令、预期现象、观察点和对应源码位置写入章节正文。这样做的好处是，实验内容不只停留在文字说明，而是可以被当前工程实际验证。

\section{项目定位}
\label{sec:intro-positioning}

DBMS\_C 是一个基于 C 语言实现的轻量化教学型 DBMS 原型，已经包含 SQL 解析、查询计划、扫描执行、记录管理、元数据管理、事务日志、缓冲区管理、文件页式存储等模块。它强调“可运行、可观察、可阅读”，适合作为数据库系统课程中的实验材料和源码阅读对象。

需要注意的是，本项目不是完整工业级数据库系统。它不会覆盖工业数据库中的全部 SQL 标准、优化器规则、并发事务隔离级别、存储引擎能力和高性能工程化细节。本指导书也不是论文式说明文档，而是配套 DBMS\_C 源码的实验指导书。

\section{学习主线}
\label{sec:intro-learning-path}

阅读 DBMS\_C 时，可以沿着从用户输入向底层存储逐步下沉的方向理解系统。图 \ref{fig:intro-learning-path} 不是完整函数调用图，而是一张学习路线图：上层更接近 SQL 和用户观察结果，下层更接近数据页、文件和事务支撑。

\begin{figure}[htbp]
  \centering
  \resizebox{0.92\textwidth}{!}{%
  \begin{tikzpicture}[node distance=6mm and 12mm]
    \node[dblevel] (sql) {SQL 输入\\\codepath{main.c}};
    \node[dbnode, below=of sql] (parse) {SQL 解析\\\codepath{parse/}};
    \node[dbnode, below=of parse] (plan) {查询计划\\\codepath{plan/}};
    \node[dbnode, below=of plan] (scan) {Scan 扫描执行\\\codepath{query/}};
    \node[dbnode, below=of scan] (record) {记录访问\\\codepath{record/}};
    \node[dbnode, below left=8mm and 10mm of record] (metadata) {元数据支撑\\\codepath{metadata/}};
    \node[dbnode, below=of record] (tx) {事务与日志\\\codepath{tx/} + \codepath{Log/}};
    \node[dbnode, below=of tx] (buffer) {缓冲区管理\\\codepath{buffer/}};
    \node[dblevel, below=of buffer] (file) {文件与页式存储\\\codepath{File/}};
    \node[dbsubnode, right=22mm of plan] (trace) {Trace 可视化观察\\\codepath{trace/}};

    \draw[dbdown] (sql) -- (parse);
    \draw[dbdown] (parse) -- (plan);
    \draw[dbdown] (plan) -- (scan);
    \draw[dbdown] (scan) -- (record);
    \draw[dbdown] (record) -- (tx);
    \draw[dbdown] (tx) -- (buffer);
    \draw[dbdown] (buffer) -- (file);
    \draw[dbarrow] (metadata) -- (record);
    \draw[dbarrow] (parse.east) -- ++(10mm,0) |- (trace.west);
    \draw[dbarrow] (plan.east) -- (trace.west);
    \draw[dbarrow] (scan.east) -- ++(10mm,0) |- (trace.west);

    \begin{scope}[on background layer]
      \node[dbgroup, fit=(sql)(parse)(plan)(scan)(record)(tx)(buffer)(file), label={[font=\small\color{DBMSGray}]left:层次下沉}] {};
    \end{scope}
  \end{tikzpicture}}
  \caption{DBMS\_C 学习主线图}
  \label{fig:intro-learning-path}
\end{figure}

图 \ref{fig:intro-learning-path} 中，各层的作用可以这样理解：SQL 输入层负责接收用户命令；SQL 解析层把文本转换为内部数据对象；查询计划层决定如何组合表扫描、选择和投影；Scan 扫描执行层以统一接口向上提供记录；记录访问层负责表记录的定位与读写；元数据层保存表结构、视图和索引相关信息；事务与日志层为更新和恢复提供支撑；缓冲区层管理内存中的页；文件与页式存储层最终负责磁盘文件和页面读写。Trace 模式不改变执行语义，只帮助读者观察这些阶段的大致过程。

\section{项目目录概览}
\label{sec:intro-directories}

DBMS\_C 的目录不是按论文章节划分，而是按数据库执行链路和模块职责组织。表 \ref{tab:intro-directories} 根据当前仓库中的真实目录整理，便于读者在各章实验中快速定位源码。

\begin{table}[htbp]
  \centering
  \caption{DBMS\_C 项目目录职责表}
  \label{tab:intro-directories}
  \begin{tabularx}{\textwidth}{>{\ttfamily}l Y}
    \toprule
    目录或文件 & 主要职责 \\
    \midrule
    Lib/ & 基础工具库，包括字符串、字节缓冲、列表、映射、向量、分词辅助等通用组件。 \\
    File/ & 文件块、页对象和文件管理，是理解页式存储的入口。 \\
    Log/ & 日志管理相关实现，为事务提交、回滚和恢复提供基础。 \\
    buffer/ & 缓冲区管理和 LRU 替换策略，负责内存页与磁盘页之间的协调。 \\
    tx/ & 事务、并发控制、恢复管理以及事务持有的缓冲区列表。 \\
    record/ & Schema、Layout、RID、RecordPage、TableScan 等记录层结构。 \\
    metadata/ & 表、字段、视图、索引和统计信息等元数据管理。 \\
    parse/ & Lexer、Parser 以及各类 SQL 命令的数据对象。 \\
    plan/ & 查询计划和更新计划的构造入口。 \\
    query/ & Scan 接口、表达式、谓词、选择、投影和笛卡尔积扫描等执行层对象。 \\
    index/ & 当前仓库中的哈希索引相关实现。 \\
    DBMS/ & SimpleDB 初始化入口，连接文件、日志、缓冲区、元数据和规划器。 \\
    main.c & CLI 入口，负责读取命令、分发 SQL、执行事务命令和输出结果。 \\
    test/ & 基于 CMocka 的已有测试文件。 \\
    docs/ & 说明文档、Trace 文档和本实验指导书。 \\
    \bottomrule
  \end{tabularx}
\end{table}

推荐阅读顺序是：\codepath{Lib/} → \codepath{File/} → \codepath{Log/} → \codepath{buffer/} → \codepath{tx/} → \codepath{record/} → \codepath{metadata/} → \codepath{parse/} → \codepath{plan/} 与 \codepath{query/} → 更新执行 → Trace → 综合实验。这个顺序先建立底层对象概念，再理解上层 SQL 如何使用这些能力。

\section{环境准备}
\label{sec:intro-environment}

本指导书默认在 Windows 环境下进行实验。当前项目的稳定构建方式依赖 CMake、MinGW 编译器和 \codepath{mingw32-make}。编译本指导书本身还需要 XeLaTeX。

建议准备以下环境：

\begin{itemize}
  \item Windows 操作系统；
  \item MinGW 或等价的 GCC 工具链，并确保 \codepath{mingw32-make} 可用；
  \item CMake；
  \item XeLaTeX，用于编译本指导书；
  \item 可选编辑器：VS Code、CLion、Trae 或 Codex。
\end{itemize}

如果不同机器上的工具链路径不同，应优先保证命令行中能直接执行 \codepath{cmake}、\codepath{mingw32-make} 和 \codepath{xelatex}。本指导书不会假定某一个 IDE 的专有功能。

\section{构建项目}
\label{sec:intro-build}

当前 Windows/Codex 环境下，DBMS\_C 已验证的稳定构建命令如表 \ref{tab:intro-build-test-commands} 所示。请先进入项目根目录，再进行配置和编译。这里的“项目根目录”指包含 \codepath{CMakeLists.txt} 的目录；如果读者的项目位于其他位置，应替换为自己本机的实际路径。例如项目可以位于 \path{D:/code/DBMS/NewDBMS/DBMS_C}，但本指导书后续命令统一假设已经进入项目根目录。

\begin{table}[htbp]
  \centering
  \caption{构建与测试命令表}
  \label{tab:intro-build-test-commands}
  \begin{tabularx}{\textwidth}{>{\ttfamily}l Y}
    \toprule
    命令 & 作用 \\
    \midrule
    cd <project-root> & 进入包含 \codepath{CMakeLists.txt} 的项目根目录。 \\
    cmake -S . -B build & 根据当前 \codepath{CMakeLists.txt} 配置构建目录。 \\
    mingw32-make -C build -j8 SHELL=cmd.exe & 使用 \codepath{mingw32-make} 编译项目，并显式覆盖 Makefile 中的 shell。 \\
    ctest --test-dir build --output-on-failure & 运行已注册测试，并在失败时输出详细信息。 \\
    bin/NewDBMS.exe & 启动交互式 DBMS 程序。 \\
    \bottomrule
  \end{tabularx}
\end{table}

\begin{dbcode}[listing options={language=bash,numbers=none}]{稳定构建命令}
cmake -S . -B build
mingw32-make -C build -j8 SHELL=cmd.exe
\end{dbcode}

这里使用 \codepath{mingw32-make -C build -j8 SHELL=cmd.exe}，是因为当前构建目录在 Windows 下使用 Makefile 时，默认 shell 可能指向 \codepath{/bin/sh}。在 Codex/Trae 这类环境中，调用 Git/MSYS 的 shell 曾经触发过构建阶段问题。因此，本指导书把 \codepath{mingw32-make} 并覆盖 \codepath{SHELL=cmd.exe} 作为当前稳定主命令，不把 \codepath{cmake --build build} 作为默认主命令。

\section{运行测试与测试驱动的实验组织方式}
\label{sec:intro-test-driven}

项目测试通过 CTest 统一运行：

\begin{dbcode}[listing options={language=bash,numbers=none}]{运行已有测试}
ctest --test-dir build --output-on-failure
\end{dbcode}

\codepath{ctest} 用于确认已有模块测试是否通过。每章新增或调整测试后，都要重新运行完整 CTest，确保没有破坏已有测试。如果新增测试，测试数量可能增加；如果测试失败，应优先查看第一条真实失败信息，而不是同时修改多个无关位置。

从第1章开始，本指导书采用测试驱动的实验组织方式：

\begin{enumerate}
  \item 每个模块实验均以真实可运行的测试或明确的人工验证方式作为章节依据。
  \item 测试必须基于当前仓库真实 API，不能编造不存在的函数、结构体或路径。
  \item 测试通过后，章节中再写入测试文件路径、测试目标、运行命令、预期结果、观察点和对应源码位置。
  \item 文档中不能写不存在或没有跑通的测试。
  \item 如果未来扩展新的实验主题，应先建立真实验证方式，不能伪造测试结果。
  \item 每新增一个测试后都要运行完整 CTest。
  \item 第0章本身不新增模块测试，只记录当前工程的基线构建与测试方法。
\end{enumerate}

表 \ref{tab:intro-chapter-test-strategy} 汇总了本指导书各章节与配套测试之间的对应关系。除第11章 Trace 可视化实验采用人工验证方式外，其余核心模块均配套了可运行的 CTest 测试，用于支撑章节中的源码讲解和实验观察。

\begin{table}[htbp]
  \centering
  \caption{章节与测试策略对照表}
  \label{tab:intro-chapter-test-strategy}
  \begin{tabularx}{\textwidth}{>{\raggedright\arraybackslash}p{0.25\textwidth} Y}
    \toprule
    章节 & 配套测试与验证内容 \\
    \midrule
    第1章 基础组件与工具库实验 & LibBasicTest：验证 CString、CVector、CList、CMap、ByteBuffer、StreamTokenizer 等基础组件。 \\
    第2章 文件与页式存储实验 & FileManagerBasicTest：验证 BlockID、Page、FileManager 的页读写与追加。 \\
    第3章 日志管理实验 & LogManagerBasicTest：验证日志追加、刷盘、逆序读取和重新打开读回。 \\
    第4章 缓冲区管理实验 & BufferManagerBasicTest：验证 pin/unpin、重复 pin、脏页 flush 和 frame 复用。 \\
    第5章 事务、并发与恢复实验 & TransactionBasicTest、ConcurrencyBasicTest、RecoveryBasicTest：验证事务基础、S/X 锁和 rollback 回滚。 \\
    第6章 记录管理实验 & RecordBasicTest：验证 Schema、Layout、RID、RecordPage 和 TableScan。 \\
    第7章 元数据管理实验 & MetadataBasicTest：验证表、视图、索引元数据登记与读取。 \\
    第8章 SQL 解析实验 & ParseBasicTest：验证 SELECT、INSERT、UPDATE、DELETE、CREATE TABLE、CREATE VIEW、CREATE INDEX 的解析结果。 \\
    第9章 查询计划与扫描执行实验 & PlanScanBasicTest：验证 QueryData 到 Plan，再到 Scan 的查询执行路径。 \\
    第10章 更新执行实验 & UpdateBasicTest：验证 CREATE TABLE、INSERT、UPDATE、DELETE、CREATE VIEW、CREATE INDEX 的更新执行路径。 \\
    第11章 SQL 执行过程 Trace 可视化实验 & 人工验证：通过 CLI Trace 观察 Parse、Plan、Scan、Update、Transaction 等阶段输出。 \\
    第12章 综合实验与总结 & 综合说明：串联全书模块、测试体系、运行方式和扩展方向。 \\
    \bottomrule
  \end{tabularx}
\end{table}

\section{运行系统}
\label{sec:intro-run-system}

构建完成后，主程序位于项目根目录下的 \codepath{bin/NewDBMS.exe}。在 Windows PowerShell 中可以这样启动：

\begin{dbcode}[listing options={language=bash,numbers=none}]{启动 DBMS\_C}
./bin/NewDBMS.exe
\end{dbcode}

程序启动后会进入交互式 SQL 提示符。读者可以输入 SQL 语句并以分号结束，也可以输入 \codepath{commit}、\codepath{rollback}、\codepath{trace on}、\codepath{trace off}、\codepath{exit} 等元命令。输入 \codepath{exit} 会退出系统；按当前程序逻辑，退出阶段会提交当前事务。因此，如果 Trace 已开启，读者可能在关闭数据库时看到一段退出阶段的 \codepath{COMMIT} Trace。

\section{Trace 模式说明}
\label{sec:intro-trace}

Trace 模式用于观察 SQL 在 DBMS\_C 内部的大致执行过程，包括解析、计划、扫描、更新、事务和输出等阶段。它的定位是 CLI 教学可视化，不是 GUI，也不改变普通 SQL 执行语义。

当前仓库中已经有 \codepath{docs/trace_mode.md}，用于记录 Trace 模式的开启方式、示例 SQL 和当前限制。本指导书第11章会进一步围绕 Trace 设计实验；第0章只要求读者知道 Trace 是一种观察执行链路的辅助工具。

\section{如何阅读本指导书}
\label{sec:intro-how-to-read}

建议按以下方式使用本指导书：

\begin{enumerate}
  \item 先阅读每章实验目的，明确本章要理解的问题。
  \item 再查看相关源码文件，确认结构体、函数和目录都来自真实仓库。
  \item 再运行本章配套测试，观察测试是否通过以及输出是否符合预期。
  \item 对照图、表、代码和 Trace 输出理解执行链路。
  \item 最后完成思考题，把模块放回 DBMS 的整体执行路径中理解。
\end{enumerate}

阅读时可以把每章的测试名称、源码路径、图表和 Trace 观察点相互对照。第11章采用人工验证方式，重点在于观察 CLI Trace 输出与前面章节测试过的核心链路之间的对应关系。

\section{本章小结}
\label{sec:intro-summary}

本章完成了指导书的起步准备：明确了 DBMS\_C 的教学型原型定位，给出了从 SQL 输入向底层存储下沉的学习主线，整理了真实项目目录职责，说明了 Windows 环境下稳定构建、测试和运行系统的方法，并概括了全书采用测试驱动的实验组织方式。

从下一章开始，读者可以按照章节顺序阅读源码、运行对应测试，并把测试入口、源码位置和观察结果联系起来理解 DBMS\_C 的完整执行链路。
